%% kapa: exam=1 round=21 year=2010 period=1 form=A subject=economics subject_ko=경제원론 seq=4 number=44
%% source: Q-Net 감정평가사 기출문제 (cst00309) artlSeq=5024181 file=extracted/1008/5024181/2010년 제21회 감정평가사 1차시험 기출문제/1교시.hwp
%% conversion: kordoc 4.24.51 -> tools/gen1.py (auto); review=proofread-sonnet; flags=none
\begin{question}{44}
독점기업이 50단위의 재화를 생산하여 10,000원의 총수입을 얻고 있다. 이 기업의 한계비용곡선이 우상향하고 50단위의 재화를 생산할 때 한계비용은 200원이다. 다음 설명 중 옳은 것은?
\begin{choices}
  \item 생산량을 감소시킴으로써 이윤을 증가시킬 수 있다.
  \item 생산량을 증가시킴으로써 이윤을 증가시킬 수 있다.
  \item 가격을 인하시킴으로써 이윤을 증가시킬 수 있다.
  \item 이윤극대화가 달성되는 산출량을 생산하고 있다.
  \item 이윤극대화가 달성되는 가격을 설정하고 있다.
\end{choices}
\end{question}
\answer{1}
